% ----------------------------------------------------------------------
%  1.7. Кристаллография границ зёрен
% ----------------------------------------------------------------------
\begin{orientir}
Граница зерна~— место встречи двух одинаковых решёток с разной
ориентацией. Сначала выясним, \vazhno{сколько параметров} задают границу и
чем границы наклона отличаются от границ кручения. Затем разберём
\vazhno{малоугловые} границы~— стенки и сетки дислокаций (1.7.1) и
\vazhno{высокоугловые}, строение которых описывают решётки совпадающих
узлов (1.7.2).
\end{orientir}

Во всяком поликристалле есть внутренние границы: между зёрнами одной фазы
(\term{границы зёрен}) и между разными фазами (\term{межфазные границы}).
Граница~— двумерный дефект: в двух измерениях её размеры макроскопические,
а в третьем (толщина)~— атомные. Поэтому граница~— естественный
\emph{наноструктурный} элемент материала.

\begin{masshtab}[граница зерна]{-9.5/-9/{толщина 0,5--1 нм}, -8/-7.4/{дислокации в субгранице}, -5.5/-3.7/{размер зерна}}
Граница тонкая~— несколько межатомных расстояний (0,5--1~нм), но
протяжённая~— во весь размер зерна (1--100~мкм). В малоугловой
субгранице дислокации отстоят на $D=b/\theta$: при $\theta=1^\circ$ это
около 15~нм. Суммарная площадь границ огромна: в 1~см$^3$ металла с
зерном 50~мкм она около $2/d\approx400$~см$^2$. По этой «внутренней
поверхности» быстрее идёт диффузия и на ней скапливаются примеси.
\end{masshtab}

\paragraph{Сколько параметров задают границу}
Рассмотрим простейший случай~— плоскую границу в бикристалле. С
макроскопической точки зрения нужны \vazhno{5 параметров}:
\begin{itemize}
\item 3~— для взаимного разворота кристаллов: три угла Эйлера
$(\varphi,\theta,\psi)$ или угол $\theta$ и ось разворота $\vb u\parallel[mnp]$
(у единичного вектора оси независимы лишь две компоненты, так как
$l_1^2+l_2^2+l_3^2=1$);
\item 2~— для нормали $\vb n$ к плоскости границы.
\end{itemize}
На атомном уровне добавляется вектор $\vb t=t_1\vb a_1+t_2\vb a_2+t_3\vb a_3$
относительного сдвига решёток (3 параметра в долях трансляций), т.\,е.
всего \vazhno{8 параметров}.

\begin{opredelenie}
Если ось разворота лежит в плоскости границы ($\vb u\perp\vb n$), граница
называется \term{границей наклона}; если ось перпендикулярна границе
($\vb u\parallel\vb n$)~— \term{границей кручения} (рис.~\ref{fig:tilt-twist}).
Границы с углом разориентации $\theta\lesssim10$--$15^\circ$ называют
\term{малоугловыми} (это границы субзёрен), с бо́льшим~—
\term{высокоугловыми}.
\end{opredelenie}

\begin{figure}[ht]
\centering
\begin{tikzpicture}[cubeproj, scale=1.1]
  % (а) граница наклона: граница — плоскость x = 0, ось u — вдоль y
  \begin{scope}
    \foreach \x/\c in {-2/fmBlueBg, 0/fmRedBg}{
      \fill[\c] (\x,0,0) -- (\x+2,0,0) -- (\x+2,0,2) -- (\x,0,2) -- cycle;
      \fill[\c!70!white] (\x,0,2) -- (\x+2,0,2) -- (\x+2,1.6,2) -- (\x,1.6,2) -- cycle;}
    \fill[fmRedBg!60!white] (2,0,0) -- (2,1.6,0) -- (2,1.6,2) -- (2,0,2) -- cycle;
    \begin{scope}
      \clip (-2,0,0) -- (0,0,0) -- (0,0,2) -- (-2,0,2) -- cycle;
      \foreach \k in {0,...,16}{\draw[fmBlue!70, line width=0.5pt] ({-3+0.3*\k},0,-0.5) -- ++({3*cos(78)},0,{3*sin(78)});}
    \end{scope}
    \begin{scope}
      \clip (0,0,0) -- (2,0,0) -- (2,0,2) -- (0,0,2) -- cycle;
      \foreach \k in {0,...,16}{\draw[fmRed!70, line width=0.5pt] ({-0.8+0.3*\k},0,-0.5) -- ++({3*cos(102)},0,{3*sin(102)});}
    \end{scope}
    \draw[atomLine] (-2,0,0) -- (2,0,0) -- (2,0,2) -- (-2,0,2) -- cycle
      (-2,0,2) -- (-2,1.6,2) -- (2,1.6,2) -- (2,0,2) (2,0,0) -- (2,1.6,0) -- (2,1.6,2);
    \draw[fmGray, line width=1.3pt] (0,0,0) -- (0,0,2) -- (0,1.6,2);
    \draw[bvec, fmGreen!70!black] (0,-0.2,2.3) -- (0,1.9,2.3) node[above right=-2pt, font=\small]{$\vb u$};
    \draw[bvec, fmBlue] (0,0,2.6) -- (0.9,0,2.6) node[right, font=\small]{$\vb n$};
    \node[font=\small] at (0,0,-0.5) {\textit{а}};
  \end{scope}
  % (б) граница кручения: граница — плоскость z = 1, ось u — вдоль z (раздвинуто)
  \begin{scope}[shift={(4.6,0,0)}]
    \foreach \z/\c/\a in {0/fmBlueBg/-10, 1.6/fmRedBg/10}{
      \fill[\c] (0,0,\z) -- (2,0,\z) -- (2,0,\z+1) -- (0,0,\z+1) -- cycle;
      \fill[\c!60!white] (2,0,\z) -- (2,2,\z) -- (2,2,\z+1) -- (2,0,\z+1) -- cycle;
      \fill[\c!80!white] (0,0,\z+1) -- (2,0,\z+1) -- (2,2,\z+1) -- (0,2,\z+1) -- cycle;
      \begin{scope}
        \clip (0,0,\z+1) -- (2,0,\z+1) -- (2,2,\z+1) -- (0,2,\z+1) -- cycle;
        \foreach \k in {-6,...,6}{
          \draw[atomLine!60, line width=0.4pt] ({1+0.3*\k*cos(\a)-2*sin(\a)},{1+0.3*\k*sin(\a)+2*cos(\a)},\z+1)
            -- ({1+0.3*\k*cos(\a)+2*sin(\a)},{1+0.3*\k*sin(\a)-2*cos(\a)},\z+1);}
      \end{scope}
      \draw[atomLine] (0,0,\z) -- (2,0,\z) -- (2,0,\z+1) -- (0,0,\z+1) -- cycle
        (2,0,\z) -- (2,2,\z) -- (2,2,\z+1) -- (0,2,\z+1) -- (0,0,\z+1) (2,0,\z+1) -- (2,2,\z+1);}
    \draw[bvec, fmGreen!70!black] (1,1,2.6) -- (1,1,3.6) node[above, font=\small]{$\vb u\parallel\vb n$};
    \draw[fmGreen!70!black, dashed] (1,1,1.0) -- (1,1,2.6);
    \node[font=\footnotesize, anchor=west] at (2.15,0,1.3) {граница};
    \draw[fmGray, densely dotted] (2,0,1.0) -- (2,0,1.6);
    \node[font=\small] at (1,0,-0.5) {\textit{б}};
  \end{scope}
\end{tikzpicture}
\caption{Граница наклона (\textit{а}): ось разворота $\vb u$ лежит в плоскости
границы, атомные плоскости двух зёрен наклонены в разные стороны.
Граница кручения (\textit{б}, зёрна раздвинуты): ось разворота
перпендикулярна границе, и решётки двух зёрен повёрнуты относительно друг
друга в плоскости границы}
\label{fig:tilt-twist}
\end{figure}

\subsection{Малоугловые границы}\label{ssec:low-angle}

Малоугловые границы~— \emph{дислокационные} структуры (подробнее~— гл.~2,
п.~2.2.5). \term{Симметричная граница наклона}~— стенка параллельных
краевых дислокаций одного знака (рис.~\ref{fig:lab-wall}). Если расстояние
между дислокациями $D$, вектор Бюргерса $b$ и угол разориентации $\theta$,
то $\sin(\theta/2)=b/2D$, а при малых углах
\begin{keyeq}\label{eq:wall-spacing}
D=\frac b\theta .
\end{keyeq}
Чем больше угол, тем ближе дислокации друг к другу. При $\theta\gtrsim10$--$15^\circ$
расстояние $D$ становится порядка нескольких $b$, ядра дислокаций
сливаются, и дислокационная модель теряет смысл~— граница становится
высокоугловой.

\begin{figure}[ht]
\centering
\begin{tikzpicture}[scale=0.42]
  % Сгенерировано tools/gen_tilt_boundary.py — не редактировать вручную
\begin{scope}
\clip (-6.30,-6.30) rectangle (6.30,6.30);
\draw[atomLine!55, line width=0.45pt] (-7.059,-0.102)--(-7.178,0.897) (-7.059,-0.102)--(-6.940,-1.008) (-7.059,-0.102)--(-6.060,-0.008) (-7.178,0.897)--(-7.297,1.896) (-7.178,0.897)--(-6.180,0.992) (-7.297,1.896)--(-6.299,1.990) (-6.535,3.893)--(-6.653,4.890) (-6.535,3.893)--(-6.419,2.989) (-6.535,3.893)--(-5.539,3.987) (-6.653,4.890)--(-6.772,5.887) (-6.653,4.890)--(-5.658,4.984) (-6.772,5.887)--(-6.890,6.884) (-6.772,5.887)--(-5.778,5.980) (-6.890,6.884)--(-5.897,6.976) (-7.222,-7.020)--(-6.216,-6.922) (-6.581,-4.011)--(-6.701,-3.010) (-6.581,-4.011)--(-6.457,-4.916) (-6.581,-4.011)--(-5.578,-3.913) (-6.701,-3.010)--(-6.820,-2.008) (-6.701,-3.010)--(-5.699,-2.911) (-6.820,-2.008)--(-6.940,-1.008) (-6.820,-2.008)--(-5.819,-1.910) (-6.940,-1.008)--(-5.940,-0.910) (-6.060,-0.008)--(-6.180,0.992) (-6.060,-0.008)--(-5.940,-0.910) (-6.060,-0.008)--(-5.060,0.090) (-6.180,0.992)--(-6.299,1.990) (-6.180,0.992)--(-5.181,1.090) (-6.299,1.990)--(-6.419,2.989) (-6.299,1.990)--(-5.301,2.089) (-6.419,2.989)--(-5.422,3.087) (-5.539,3.987)--(-5.658,4.984) (-5.539,3.987)--(-5.422,3.087) (-5.539,3.987)--(-4.543,4.084) (-5.658,4.984)--(-5.778,5.980) (-5.658,4.984)--(-4.663,5.081) (-5.778,5.980)--(-5.897,6.976) (-5.778,5.980)--(-4.784,6.078) (-5.897,6.976)--(-4.904,7.073) (-6.216,-6.922)--(-6.337,-5.919) (-6.216,-6.922)--(-5.211,-6.820) (-6.337,-5.919)--(-6.457,-4.916) (-6.337,-5.919)--(-5.331,-5.816) (-6.457,-4.916)--(-5.453,-4.814) (-5.578,-3.913)--(-5.699,-2.911) (-5.578,-3.913)--(-5.453,-4.814) (-5.578,-3.913)--(-4.575,-3.811) (-5.699,-2.911)--(-5.819,-1.910) (-5.699,-2.911)--(-4.697,-2.808) (-5.819,-1.910)--(-5.940,-0.910) (-5.819,-1.910)--(-4.817,-1.807) (-5.940,-0.910)--(-4.939,-0.807) (-5.060,0.090)--(-5.181,1.090) (-5.060,0.090)--(-4.939,-0.807) (-5.060,0.090)--(-4.061,0.193) (-5.181,1.090)--(-5.301,2.089) (-5.181,1.090)--(-4.183,1.193) (-5.301,2.089)--(-5.422,3.087) (-5.301,2.089)--(-4.303,2.192) (-5.422,3.087)--(-4.425,3.189) (-4.543,4.084)--(-4.663,5.081) (-4.543,4.084)--(-4.425,3.189) (-4.543,4.084)--(-3.547,4.186) (-4.663,5.081)--(-4.784,6.078) (-4.663,5.081)--(-3.669,5.184) (-4.784,6.078)--(-4.904,7.073) (-4.784,6.078)--(-3.789,6.180) (-4.904,7.073)--(-3.910,7.174) (-5.211,-6.820)--(-5.331,-5.816) (-5.211,-6.820)--(-4.208,-6.713) (-5.331,-5.816)--(-5.453,-4.814) (-5.331,-5.816)--(-4.325,-5.709) (-5.453,-4.814)--(-4.447,-4.708) (-4.575,-3.811)--(-4.697,-2.808) (-4.575,-3.811)--(-4.447,-4.708) (-4.575,-3.811)--(-3.574,-3.706) (-4.697,-2.808)--(-4.817,-1.807) (-4.697,-2.808)--(-3.697,-2.700) (-4.817,-1.807)--(-4.939,-0.807) (-4.817,-1.807)--(-3.814,-1.699) (-4.939,-0.807)--(-3.937,-0.701) (-4.061,0.193)--(-4.183,1.193) (-4.061,0.193)--(-3.937,-0.701) (-4.061,0.193)--(-3.063,0.299) (-4.183,1.193)--(-4.303,2.192) (-4.183,1.193)--(-3.186,1.301) (-4.303,2.192)--(-4.425,3.189) (-4.303,2.192)--(-3.303,2.300) (-4.425,3.189)--(-3.426,3.294) (-3.547,4.186)--(-3.669,5.184) (-3.547,4.186)--(-3.426,3.294) (-3.547,4.186)--(-2.553,4.292) (-3.669,5.184)--(-3.789,6.180) (-3.669,5.184)--(-2.675,5.291) (-3.789,6.180)--(-3.910,7.174) (-3.789,6.180)--(-2.792,6.287) (-3.910,7.174)--(-2.915,7.279) (-4.208,-6.713)--(-4.325,-5.709) (-4.208,-6.713)--(-3.208,-6.598) (-4.325,-5.709)--(-4.447,-4.708) (-4.325,-5.709)--(-3.315,-5.594) (-4.447,-4.708)--(-3.438,-4.601) (-3.574,-3.706)--(-3.697,-2.700) (-3.574,-3.706)--(-3.438,-4.601) (-3.574,-3.706)--(-2.578,-3.598) (-3.697,-2.700)--(-3.814,-1.699) (-3.697,-2.700)--(-2.701,-2.585) (-3.814,-1.699)--(-3.937,-0.701) (-3.814,-1.699)--(-2.807,-1.583) (-3.937,-0.701)--(-2.930,-0.594) (-3.063,0.299)--(-3.186,1.301) (-3.063,0.299)--(-2.930,-0.594) (-3.063,0.299)--(-2.070,0.406) (-3.186,1.301)--(-3.303,2.300) (-3.186,1.301)--(-2.193,1.417) (-3.303,2.300)--(-3.426,3.294) (-3.303,2.300)--(-2.299,2.415) (-3.426,3.294)--(-2.422,3.402) (-2.553,4.292)--(-2.675,5.291) (-2.553,4.292)--(-2.422,3.402) (-2.553,4.292)--(-1.562,4.399) (-2.675,5.291)--(-2.792,6.287) (-2.675,5.291)--(-1.685,5.406) (-2.792,6.287)--(-2.915,7.279) (-2.792,6.287)--(-1.792,6.402) (-3.208,-6.598)--(-3.315,-5.594) (-3.208,-6.598)--(-2.221,-6.473) (-3.315,-5.594)--(-3.438,-4.601) (-3.315,-5.594)--(-2.293,-5.469) (-3.438,-4.601)--(-2.411,-4.495) (-2.578,-3.598)--(-2.701,-2.585) (-2.578,-3.598)--(-2.411,-4.495) (-2.578,-3.598)--(-1.598,-3.492) (-2.701,-2.585)--(-2.807,-1.583) (-2.701,-2.585)--(-1.716,-2.460) (-2.807,-1.583)--(-2.930,-0.594) (-2.807,-1.583)--(-1.789,-1.458) (-2.930,-0.594)--(-1.906,-0.487) (-2.070,0.406)--(-2.193,1.417) (-2.070,0.406)--(-1.906,-0.487) (-2.070,0.406)--(-1.094,0.513) (-2.193,1.417)--(-2.299,2.415) (-2.193,1.417)--(-1.211,1.542) (-2.299,2.415)--(-2.422,3.402) (-2.299,2.415)--(-1.284,2.540) (-2.422,3.402)--(-1.402,3.508) (-1.562,4.399)--(-1.685,5.406) (-1.562,4.399)--(-1.402,3.508) (-1.562,4.399)--(-0.589,4.505) (-1.685,5.406)--(-1.792,6.402) (-1.685,5.406)--(-0.707,5.531) (-1.792,6.402)--(-0.779,6.527) (-2.221,-6.473)--(-2.293,-5.469) (-2.221,-6.473)--(-1.244,-6.356) (-2.293,-5.469)--(-2.411,-4.495) (-2.293,-5.469)--(-1.261,-5.352) (-2.411,-4.495)--(-1.317,-4.365) (-1.598,-3.492)--(-1.716,-2.460) (-1.598,-3.492)--(-1.317,-4.365) (-1.598,-3.492)--(-0.686,-3.362) (-1.716,-2.460)--(-1.789,-1.458) (-1.716,-2.460)--(-0.742,-2.343) (-1.789,-1.458)--(-1.906,-0.487) (-1.789,-1.458)--(-0.759,-1.341) (-1.906,-0.487)--(-0.816,-0.357) (-1.094,0.513)--(-1.211,1.542) (-1.094,0.513)--(-0.816,-0.357) (-1.094,0.513)--(-0.184,0.643) (-1.211,1.542)--(-1.284,2.540) (-1.211,1.542)--(-0.241,1.659) (-1.284,2.540)--(-1.402,3.508) (-1.284,2.540)--(-0.258,2.657) (-1.402,3.508)--(-0.314,3.638) (-0.589,4.505)--(-0.707,5.531) (-0.589,4.505)--(-0.314,3.638) (-0.589,4.505)--(0.317,4.635) (-0.707,5.531)--(-0.779,6.527) (-0.707,5.531)--(0.261,5.648) (-0.779,6.527)--(0.244,6.644) (-1.244,-6.356)--(-1.261,-5.352) (-1.244,-6.356)--(-0.256,-6.356) (-1.261,-5.352)--(-1.317,-4.365) (-1.261,-5.352)--(-0.239,-5.352) (-1.317,-4.365)--(-0.686,-3.362) (-1.317,-4.365)--(-0.183,-4.365) (-0.686,-3.362)--(-0.742,-2.343) (-0.686,-3.362)--(-0.183,-4.365) (-0.686,-3.362)--(0.186,-3.362) (-0.742,-2.343)--(-0.759,-1.341) (-0.742,-2.343)--(0.242,-2.343) (-0.759,-1.341)--(-0.816,-0.357) (-0.759,-1.341)--(0.259,-1.341) (-0.816,-0.357)--(-0.184,0.643) (-0.816,-0.357)--(0.316,-0.357) (-0.184,0.643)--(-0.241,1.659) (-0.184,0.643)--(0.316,-0.357) (-0.184,0.643)--(0.684,0.643) (-0.241,1.659)--(-0.258,2.657) (-0.241,1.659)--(0.741,1.659) (-0.258,2.657)--(-0.314,3.638) (-0.258,2.657)--(0.758,2.657) (-0.314,3.638)--(0.317,4.635) (-0.314,3.638)--(0.814,3.638) (0.317,4.635)--(0.261,5.648) (0.317,4.635)--(0.814,3.638) (0.317,4.635)--(1.183,4.635) (0.261,5.648)--(0.244,6.644) (0.261,5.648)--(1.239,5.648) (0.244,6.644)--(1.256,6.644) (-0.256,-6.356)--(-0.239,-5.352) (-0.256,-6.356)--(0.721,-6.473) (-0.239,-5.352)--(-0.183,-4.365) (-0.239,-5.352)--(0.793,-5.469) (-0.183,-4.365)--(0.186,-3.362) (-0.183,-4.365)--(0.911,-4.495) (0.186,-3.362)--(0.242,-2.343) (0.186,-3.362)--(1.098,-3.492) (0.242,-2.343)--(0.259,-1.341) (0.242,-2.343)--(1.216,-2.460) (0.259,-1.341)--(0.316,-0.357) (0.259,-1.341)--(1.289,-1.458) (0.316,-0.357)--(0.684,0.643) (0.316,-0.357)--(1.406,-0.487) (0.684,0.643)--(0.741,1.659) (0.684,0.643)--(1.594,0.513) (0.741,1.659)--(0.758,2.657) (0.741,1.659)--(1.711,1.542) (0.758,2.657)--(0.814,3.638) (0.758,2.657)--(1.784,2.540) (0.814,3.638)--(1.183,4.635) (0.814,3.638)--(1.902,3.508) (1.183,4.635)--(1.239,5.648) (1.183,4.635)--(2.089,4.505) (1.239,5.648)--(1.256,6.644) (1.239,5.648)--(2.207,5.531) (1.256,6.644)--(2.279,6.527) (0.721,-6.473)--(0.793,-5.469) (0.721,-6.473)--(1.708,-6.598) (0.793,-5.469)--(0.911,-4.495) (0.793,-5.469)--(1.815,-5.594) (0.911,-4.495)--(1.098,-3.492) (0.911,-4.495)--(1.938,-4.601) (1.098,-3.492)--(1.216,-2.460) (1.098,-3.492)--(2.078,-3.598) (1.216,-2.460)--(1.289,-1.458) (1.216,-2.460)--(2.201,-2.585) (1.289,-1.458)--(1.406,-0.487) (1.289,-1.458)--(2.307,-1.583) (1.406,-0.487)--(1.594,0.513) (1.406,-0.487)--(2.430,-0.594) (1.594,0.513)--(1.711,1.542) (1.594,0.513)--(2.570,0.406) (1.711,1.542)--(1.784,2.540) (1.711,1.542)--(2.693,1.417) (1.784,2.540)--(1.902,3.508) (1.784,2.540)--(2.799,2.415) (1.902,3.508)--(2.089,4.505) (1.902,3.508)--(2.922,3.402) (2.089,4.505)--(2.207,5.531) (2.089,4.505)--(3.062,4.399) (2.207,5.531)--(2.279,6.527) (2.207,5.531)--(3.185,5.406) (2.279,6.527)--(3.292,6.402) (1.708,-6.598)--(1.815,-5.594) (1.708,-6.598)--(2.708,-6.713) (1.815,-5.594)--(1.938,-4.601) (1.815,-5.594)--(2.825,-5.709) (1.938,-4.601)--(2.078,-3.598) (1.938,-4.601)--(2.947,-4.708) (2.078,-3.598)--(2.201,-2.585) (2.078,-3.598)--(3.074,-3.706) (2.201,-2.585)--(2.307,-1.583) (2.201,-2.585)--(3.197,-2.700) (2.307,-1.583)--(2.430,-0.594) (2.307,-1.583)--(3.314,-1.699) (2.430,-0.594)--(2.570,0.406) (2.430,-0.594)--(3.437,-0.701) (2.570,0.406)--(2.693,1.417) (2.570,0.406)--(3.563,0.299) (2.693,1.417)--(2.799,2.415) (2.693,1.417)--(3.686,1.301) (2.799,2.415)--(2.922,3.402) (2.799,2.415)--(3.803,2.300) (2.922,3.402)--(3.062,4.399) (2.922,3.402)--(3.926,3.294) (3.062,4.399)--(3.185,5.406) (3.062,4.399)--(4.053,4.292) (3.185,5.406)--(3.292,6.402) (3.185,5.406)--(4.175,5.291) (3.292,6.402)--(4.292,6.287) (2.708,-6.713)--(2.825,-5.709) (2.708,-6.713)--(3.711,-6.820) (2.825,-5.709)--(2.947,-4.708) (2.825,-5.709)--(3.831,-5.816) (2.947,-4.708)--(3.074,-3.706) (2.947,-4.708)--(3.953,-4.814) (3.074,-3.706)--(3.197,-2.700) (3.074,-3.706)--(4.075,-3.811) (3.197,-2.700)--(3.314,-1.699) (3.197,-2.700)--(4.197,-2.808) (3.314,-1.699)--(3.437,-0.701) (3.314,-1.699)--(4.317,-1.807) (3.437,-0.701)--(3.563,0.299) (3.437,-0.701)--(4.439,-0.807) (3.563,0.299)--(3.686,1.301) (3.563,0.299)--(4.561,0.193) (3.686,1.301)--(3.803,2.300) (3.686,1.301)--(4.683,1.193) (3.803,2.300)--(3.926,3.294) (3.803,2.300)--(4.803,2.192) (3.926,3.294)--(4.053,4.292) (3.926,3.294)--(4.925,3.189) (4.053,4.292)--(4.175,5.291) (4.053,4.292)--(5.047,4.186) (4.175,5.291)--(4.292,6.287) (4.175,5.291)--(5.169,5.184) (4.292,6.287)--(4.415,7.279) (4.292,6.287)--(5.289,6.180) (4.415,7.279)--(5.410,7.174) (3.711,-6.820)--(3.831,-5.816) (3.711,-6.820)--(4.716,-6.922) (3.831,-5.816)--(3.953,-4.814) (3.831,-5.816)--(4.837,-5.919) (3.953,-4.814)--(4.075,-3.811) (3.953,-4.814)--(4.957,-4.916) (4.075,-3.811)--(4.197,-2.808) (4.075,-3.811)--(5.078,-3.913) (4.197,-2.808)--(4.317,-1.807) (4.197,-2.808)--(5.199,-2.911) (4.317,-1.807)--(4.439,-0.807) (4.317,-1.807)--(5.319,-1.910) (4.439,-0.807)--(4.561,0.193) (4.439,-0.807)--(5.440,-0.910) (4.561,0.193)--(4.683,1.193) (4.561,0.193)--(5.560,0.090) (4.683,1.193)--(4.803,2.192) (4.683,1.193)--(5.681,1.090) (4.803,2.192)--(4.925,3.189) (4.803,2.192)--(5.801,2.089) (4.925,3.189)--(5.047,4.186) (4.925,3.189)--(5.922,3.087) (5.047,4.186)--(5.169,5.184) (5.047,4.186)--(6.043,4.084) (5.169,5.184)--(5.289,6.180) (5.169,5.184)--(6.163,5.081) (5.289,6.180)--(5.410,7.174) (5.289,6.180)--(6.284,6.078) (5.410,7.174)--(6.404,7.073) (4.716,-6.922)--(4.837,-5.919) (4.716,-6.922)--(5.722,-7.020) (4.837,-5.919)--(4.957,-4.916) (4.837,-5.919)--(5.842,-6.016) (4.957,-4.916)--(5.078,-3.913) (4.957,-4.916)--(5.961,-5.013) (5.078,-3.913)--(5.199,-2.911) (5.078,-3.913)--(6.081,-4.011) (5.199,-2.911)--(5.319,-1.910) (5.199,-2.911)--(6.201,-3.010) (5.319,-1.910)--(5.440,-0.910) (5.319,-1.910)--(6.320,-2.008) (5.440,-0.910)--(5.560,0.090) (5.440,-0.910)--(6.440,-1.008) (5.560,0.090)--(5.681,1.090) (5.560,0.090)--(6.560,-0.008) (5.681,1.090)--(5.801,2.089) (5.681,1.090)--(6.680,0.992) (5.801,2.089)--(5.922,3.087) (5.801,2.089)--(6.799,1.990) (5.922,3.087)--(6.043,4.084) (5.922,3.087)--(6.919,2.989) (6.043,4.084)--(6.163,5.081) (6.043,4.084)--(7.039,3.987) (6.163,5.081)--(6.284,6.078) (6.163,5.081)--(7.158,4.984) (6.284,6.078)--(6.404,7.073) (6.284,6.078)--(7.278,5.980) (5.722,-7.020)--(5.842,-6.016) (5.722,-7.020)--(6.728,-7.113) (5.842,-6.016)--(5.961,-5.013) (5.842,-6.016)--(6.847,-6.110) (5.961,-5.013)--(6.081,-4.011) (5.961,-5.013)--(6.965,-5.107) (6.081,-4.011)--(6.201,-3.010) (6.081,-4.011)--(7.084,-4.105) (6.201,-3.010)--(6.320,-2.008) (6.201,-3.010)--(7.203,-3.104) (6.320,-2.008)--(6.440,-1.008) (6.440,-1.008)--(6.560,-0.008) (6.560,-0.008)--(6.680,0.992) (6.680,0.992)--(6.799,1.990) (6.799,1.990)--(6.919,2.989) (6.919,2.989)--(7.039,3.987) (7.039,3.987)--(7.158,4.984) (7.158,4.984)--(7.278,5.980) (6.728,-7.113)--(6.847,-6.110) (6.847,-6.110)--(6.965,-5.107) (6.965,-5.107)--(7.084,-4.105) (7.084,-4.105)--(7.203,-3.104);
\fill[atomA, draw=atomLine, line width=0.35pt] (-7.059,-0.102) circle (0.20);
\fill[atomA, draw=atomLine, line width=0.35pt] (-7.178,0.897) circle (0.20);
\fill[atomA, draw=atomLine, line width=0.35pt] (-7.297,1.896) circle (0.20);
\fill[atomA, draw=atomLine, line width=0.35pt] (-6.535,3.893) circle (0.20);
\fill[atomA, draw=atomLine, line width=0.35pt] (-6.653,4.890) circle (0.20);
\fill[atomA, draw=atomLine, line width=0.35pt] (-6.772,5.887) circle (0.20);
\fill[atomA, draw=atomLine, line width=0.35pt] (-6.890,6.884) circle (0.20);
\fill[atomA, draw=atomLine, line width=0.35pt] (-7.222,-7.020) circle (0.20);
\fill[atomA, draw=atomLine, line width=0.35pt] (-6.581,-4.011) circle (0.20);
\fill[atomA, draw=atomLine, line width=0.35pt] (-6.701,-3.010) circle (0.20);
\fill[atomA, draw=atomLine, line width=0.35pt] (-6.820,-2.008) circle (0.20);
\fill[atomA, draw=atomLine, line width=0.35pt] (-6.940,-1.008) circle (0.20);
\fill[atomA, draw=atomLine, line width=0.35pt] (-6.060,-0.008) circle (0.20);
\fill[atomA, draw=atomLine, line width=0.35pt] (-6.180,0.992) circle (0.20);
\fill[atomA, draw=atomLine, line width=0.35pt] (-6.299,1.990) circle (0.20);
\fill[atomA, draw=atomLine, line width=0.35pt] (-6.419,2.989) circle (0.20);
\fill[atomA, draw=atomLine, line width=0.35pt] (-5.539,3.987) circle (0.20);
\fill[atomA, draw=atomLine, line width=0.35pt] (-5.658,4.984) circle (0.20);
\fill[atomA, draw=atomLine, line width=0.35pt] (-5.778,5.980) circle (0.20);
\fill[atomA, draw=atomLine, line width=0.35pt] (-5.897,6.976) circle (0.20);
\fill[atomA, draw=atomLine, line width=0.35pt] (-6.216,-6.922) circle (0.20);
\fill[atomA, draw=atomLine, line width=0.35pt] (-6.337,-5.919) circle (0.20);
\fill[atomA, draw=atomLine, line width=0.35pt] (-6.457,-4.916) circle (0.20);
\fill[atomA, draw=atomLine, line width=0.35pt] (-5.578,-3.913) circle (0.20);
\fill[atomA, draw=atomLine, line width=0.35pt] (-5.699,-2.911) circle (0.20);
\fill[atomA, draw=atomLine, line width=0.35pt] (-5.819,-1.910) circle (0.20);
\fill[atomA, draw=atomLine, line width=0.35pt] (-5.940,-0.910) circle (0.20);
\fill[atomA, draw=atomLine, line width=0.35pt] (-5.060,0.090) circle (0.20);
\fill[atomA, draw=atomLine, line width=0.35pt] (-5.181,1.090) circle (0.20);
\fill[atomA, draw=atomLine, line width=0.35pt] (-5.301,2.089) circle (0.20);
\fill[atomA, draw=atomLine, line width=0.35pt] (-5.422,3.087) circle (0.20);
\fill[atomA, draw=atomLine, line width=0.35pt] (-4.543,4.084) circle (0.20);
\fill[atomA, draw=atomLine, line width=0.35pt] (-4.663,5.081) circle (0.20);
\fill[atomA, draw=atomLine, line width=0.35pt] (-4.784,6.078) circle (0.20);
\fill[atomA, draw=atomLine, line width=0.35pt] (-4.904,7.073) circle (0.20);
\fill[atomA, draw=atomLine, line width=0.35pt] (-5.211,-6.820) circle (0.20);
\fill[atomA, draw=atomLine, line width=0.35pt] (-5.331,-5.816) circle (0.20);
\fill[atomA, draw=atomLine, line width=0.35pt] (-5.453,-4.814) circle (0.20);
\fill[atomA, draw=atomLine, line width=0.35pt] (-4.575,-3.811) circle (0.20);
\fill[atomA, draw=atomLine, line width=0.35pt] (-4.697,-2.808) circle (0.20);
\fill[atomA, draw=atomLine, line width=0.35pt] (-4.817,-1.807) circle (0.20);
\fill[atomA, draw=atomLine, line width=0.35pt] (-4.939,-0.807) circle (0.20);
\fill[atomA, draw=atomLine, line width=0.35pt] (-4.061,0.193) circle (0.20);
\fill[atomA, draw=atomLine, line width=0.35pt] (-4.183,1.193) circle (0.20);
\fill[atomA, draw=atomLine, line width=0.35pt] (-4.303,2.192) circle (0.20);
\fill[atomA, draw=atomLine, line width=0.35pt] (-4.425,3.189) circle (0.20);
\fill[atomA, draw=atomLine, line width=0.35pt] (-3.547,4.186) circle (0.20);
\fill[atomA, draw=atomLine, line width=0.35pt] (-3.669,5.184) circle (0.20);
\fill[atomA, draw=atomLine, line width=0.35pt] (-3.789,6.180) circle (0.20);
\fill[atomA, draw=atomLine, line width=0.35pt] (-3.910,7.174) circle (0.20);
\fill[atomA, draw=atomLine, line width=0.35pt] (-4.208,-6.713) circle (0.20);
\fill[atomA, draw=atomLine, line width=0.35pt] (-4.325,-5.709) circle (0.20);
\fill[atomA, draw=atomLine, line width=0.35pt] (-4.447,-4.708) circle (0.20);
\fill[atomA, draw=atomLine, line width=0.35pt] (-3.574,-3.706) circle (0.20);
\fill[atomA, draw=atomLine, line width=0.35pt] (-3.697,-2.700) circle (0.20);
\fill[atomA, draw=atomLine, line width=0.35pt] (-3.814,-1.699) circle (0.20);
\fill[atomA, draw=atomLine, line width=0.35pt] (-3.937,-0.701) circle (0.20);
\fill[atomA, draw=atomLine, line width=0.35pt] (-3.063,0.299) circle (0.20);
\fill[atomA, draw=atomLine, line width=0.35pt] (-3.186,1.301) circle (0.20);
\fill[atomA, draw=atomLine, line width=0.35pt] (-3.303,2.300) circle (0.20);
\fill[atomA, draw=atomLine, line width=0.35pt] (-3.426,3.294) circle (0.20);
\fill[atomA, draw=atomLine, line width=0.35pt] (-2.553,4.292) circle (0.20);
\fill[atomA, draw=atomLine, line width=0.35pt] (-2.675,5.291) circle (0.20);
\fill[atomA, draw=atomLine, line width=0.35pt] (-2.792,6.287) circle (0.20);
\fill[atomA, draw=atomLine, line width=0.35pt] (-2.915,7.279) circle (0.20);
\fill[atomA, draw=atomLine, line width=0.35pt] (-3.208,-6.598) circle (0.20);
\fill[atomA, draw=atomLine, line width=0.35pt] (-3.315,-5.594) circle (0.20);
\fill[atomA, draw=atomLine, line width=0.35pt] (-3.438,-4.601) circle (0.20);
\fill[atomA, draw=atomLine, line width=0.35pt] (-2.578,-3.598) circle (0.20);
\fill[atomA, draw=atomLine, line width=0.35pt] (-2.701,-2.585) circle (0.20);
\fill[atomA, draw=atomLine, line width=0.35pt] (-2.807,-1.583) circle (0.20);
\fill[atomA, draw=atomLine, line width=0.35pt] (-2.930,-0.594) circle (0.20);
\fill[atomA, draw=atomLine, line width=0.35pt] (-2.070,0.406) circle (0.20);
\fill[atomA, draw=atomLine, line width=0.35pt] (-2.193,1.417) circle (0.20);
\fill[atomA, draw=atomLine, line width=0.35pt] (-2.299,2.415) circle (0.20);
\fill[atomA, draw=atomLine, line width=0.35pt] (-2.422,3.402) circle (0.20);
\fill[atomA, draw=atomLine, line width=0.35pt] (-1.562,4.399) circle (0.20);
\fill[atomA, draw=atomLine, line width=0.35pt] (-1.685,5.406) circle (0.20);
\fill[atomA, draw=atomLine, line width=0.35pt] (-1.792,6.402) circle (0.20);
\fill[atomA, draw=atomLine, line width=0.35pt] (-2.221,-6.473) circle (0.20);
\fill[atomA, draw=atomLine, line width=0.35pt] (-2.293,-5.469) circle (0.20);
\fill[atomA, draw=atomLine, line width=0.35pt] (-2.411,-4.495) circle (0.20);
\fill[atomA, draw=atomLine, line width=0.35pt] (-1.598,-3.492) circle (0.20);
\fill[atomA, draw=atomLine, line width=0.35pt] (-1.716,-2.460) circle (0.20);
\fill[atomA, draw=atomLine, line width=0.35pt] (-1.789,-1.458) circle (0.20);
\fill[atomA, draw=atomLine, line width=0.35pt] (-1.906,-0.487) circle (0.20);
\fill[atomA, draw=atomLine, line width=0.35pt] (-1.094,0.513) circle (0.20);
\fill[atomA, draw=atomLine, line width=0.35pt] (-1.211,1.542) circle (0.20);
\fill[atomA, draw=atomLine, line width=0.35pt] (-1.284,2.540) circle (0.20);
\fill[atomA, draw=atomLine, line width=0.35pt] (-1.402,3.508) circle (0.20);
\fill[atomA, draw=atomLine, line width=0.35pt] (-0.589,4.505) circle (0.20);
\fill[atomA, draw=atomLine, line width=0.35pt] (-0.707,5.531) circle (0.20);
\fill[atomA, draw=atomLine, line width=0.35pt] (-0.779,6.527) circle (0.20);
\fill[atomA, draw=atomLine, line width=0.35pt] (-1.244,-6.356) circle (0.20);
\fill[atomA, draw=atomLine, line width=0.35pt] (-1.261,-5.352) circle (0.20);
\fill[atomA, draw=atomLine, line width=0.35pt] (-1.317,-4.365) circle (0.20);
\fill[atomA, draw=atomLine, line width=0.35pt] (-0.686,-3.362) circle (0.20);
\fill[atomA, draw=atomLine, line width=0.35pt] (-0.742,-2.343) circle (0.20);
\fill[atomA, draw=atomLine, line width=0.35pt] (-0.759,-1.341) circle (0.20);
\fill[atomA, draw=atomLine, line width=0.35pt] (-0.816,-0.357) circle (0.20);
\fill[atomA, draw=atomLine, line width=0.35pt] (-0.184,0.643) circle (0.20);
\fill[atomA, draw=atomLine, line width=0.35pt] (-0.241,1.659) circle (0.20);
\fill[atomA, draw=atomLine, line width=0.35pt] (-0.258,2.657) circle (0.20);
\fill[atomA, draw=atomLine, line width=0.35pt] (-0.314,3.638) circle (0.20);
\fill[atomA, draw=atomLine, line width=0.35pt] (0.317,4.635) circle (0.20);
\fill[atomA, draw=atomLine, line width=0.35pt] (0.261,5.648) circle (0.20);
\fill[atomA, draw=atomLine, line width=0.35pt] (0.244,6.644) circle (0.20);
\fill[atomA, draw=atomLine, line width=0.35pt] (-0.256,-6.356) circle (0.20);
\fill[atomA, draw=atomLine, line width=0.35pt] (-0.239,-5.352) circle (0.20);
\fill[atomA, draw=atomLine, line width=0.35pt] (-0.183,-4.365) circle (0.20);
\fill[atomA, draw=atomLine, line width=0.35pt] (0.186,-3.362) circle (0.20);
\fill[atomA, draw=atomLine, line width=0.35pt] (0.242,-2.343) circle (0.20);
\fill[atomA, draw=atomLine, line width=0.35pt] (0.259,-1.341) circle (0.20);
\fill[atomA, draw=atomLine, line width=0.35pt] (0.316,-0.357) circle (0.20);
\fill[atomA, draw=atomLine, line width=0.35pt] (0.684,0.643) circle (0.20);
\fill[atomA, draw=atomLine, line width=0.35pt] (0.741,1.659) circle (0.20);
\fill[atomA, draw=atomLine, line width=0.35pt] (0.758,2.657) circle (0.20);
\fill[atomA, draw=atomLine, line width=0.35pt] (0.814,3.638) circle (0.20);
\fill[atomA, draw=atomLine, line width=0.35pt] (1.183,4.635) circle (0.20);
\fill[atomA, draw=atomLine, line width=0.35pt] (1.239,5.648) circle (0.20);
\fill[atomA, draw=atomLine, line width=0.35pt] (1.256,6.644) circle (0.20);
\fill[atomA, draw=atomLine, line width=0.35pt] (0.721,-6.473) circle (0.20);
\fill[atomA, draw=atomLine, line width=0.35pt] (0.793,-5.469) circle (0.20);
\fill[atomA, draw=atomLine, line width=0.35pt] (0.911,-4.495) circle (0.20);
\fill[atomA, draw=atomLine, line width=0.35pt] (1.098,-3.492) circle (0.20);
\fill[atomA, draw=atomLine, line width=0.35pt] (1.216,-2.460) circle (0.20);
\fill[atomA, draw=atomLine, line width=0.35pt] (1.289,-1.458) circle (0.20);
\fill[atomA, draw=atomLine, line width=0.35pt] (1.406,-0.487) circle (0.20);
\fill[atomA, draw=atomLine, line width=0.35pt] (1.594,0.513) circle (0.20);
\fill[atomA, draw=atomLine, line width=0.35pt] (1.711,1.542) circle (0.20);
\fill[atomA, draw=atomLine, line width=0.35pt] (1.784,2.540) circle (0.20);
\fill[atomA, draw=atomLine, line width=0.35pt] (1.902,3.508) circle (0.20);
\fill[atomA, draw=atomLine, line width=0.35pt] (2.089,4.505) circle (0.20);
\fill[atomA, draw=atomLine, line width=0.35pt] (2.207,5.531) circle (0.20);
\fill[atomA, draw=atomLine, line width=0.35pt] (2.279,6.527) circle (0.20);
\fill[atomA, draw=atomLine, line width=0.35pt] (1.708,-6.598) circle (0.20);
\fill[atomA, draw=atomLine, line width=0.35pt] (1.815,-5.594) circle (0.20);
\fill[atomA, draw=atomLine, line width=0.35pt] (1.938,-4.601) circle (0.20);
\fill[atomA, draw=atomLine, line width=0.35pt] (2.078,-3.598) circle (0.20);
\fill[atomA, draw=atomLine, line width=0.35pt] (2.201,-2.585) circle (0.20);
\fill[atomA, draw=atomLine, line width=0.35pt] (2.307,-1.583) circle (0.20);
\fill[atomA, draw=atomLine, line width=0.35pt] (2.430,-0.594) circle (0.20);
\fill[atomA, draw=atomLine, line width=0.35pt] (2.570,0.406) circle (0.20);
\fill[atomA, draw=atomLine, line width=0.35pt] (2.693,1.417) circle (0.20);
\fill[atomA, draw=atomLine, line width=0.35pt] (2.799,2.415) circle (0.20);
\fill[atomA, draw=atomLine, line width=0.35pt] (2.922,3.402) circle (0.20);
\fill[atomA, draw=atomLine, line width=0.35pt] (3.062,4.399) circle (0.20);
\fill[atomA, draw=atomLine, line width=0.35pt] (3.185,5.406) circle (0.20);
\fill[atomA, draw=atomLine, line width=0.35pt] (3.292,6.402) circle (0.20);
\fill[atomA, draw=atomLine, line width=0.35pt] (2.708,-6.713) circle (0.20);
\fill[atomA, draw=atomLine, line width=0.35pt] (2.825,-5.709) circle (0.20);
\fill[atomA, draw=atomLine, line width=0.35pt] (2.947,-4.708) circle (0.20);
\fill[atomA, draw=atomLine, line width=0.35pt] (3.074,-3.706) circle (0.20);
\fill[atomA, draw=atomLine, line width=0.35pt] (3.197,-2.700) circle (0.20);
\fill[atomA, draw=atomLine, line width=0.35pt] (3.314,-1.699) circle (0.20);
\fill[atomA, draw=atomLine, line width=0.35pt] (3.437,-0.701) circle (0.20);
\fill[atomA, draw=atomLine, line width=0.35pt] (3.563,0.299) circle (0.20);
\fill[atomA, draw=atomLine, line width=0.35pt] (3.686,1.301) circle (0.20);
\fill[atomA, draw=atomLine, line width=0.35pt] (3.803,2.300) circle (0.20);
\fill[atomA, draw=atomLine, line width=0.35pt] (3.926,3.294) circle (0.20);
\fill[atomA, draw=atomLine, line width=0.35pt] (4.053,4.292) circle (0.20);
\fill[atomA, draw=atomLine, line width=0.35pt] (4.175,5.291) circle (0.20);
\fill[atomA, draw=atomLine, line width=0.35pt] (4.292,6.287) circle (0.20);
\fill[atomA, draw=atomLine, line width=0.35pt] (4.415,7.279) circle (0.20);
\fill[atomA, draw=atomLine, line width=0.35pt] (3.711,-6.820) circle (0.20);
\fill[atomA, draw=atomLine, line width=0.35pt] (3.831,-5.816) circle (0.20);
\fill[atomA, draw=atomLine, line width=0.35pt] (3.953,-4.814) circle (0.20);
\fill[atomA, draw=atomLine, line width=0.35pt] (4.075,-3.811) circle (0.20);
\fill[atomA, draw=atomLine, line width=0.35pt] (4.197,-2.808) circle (0.20);
\fill[atomA, draw=atomLine, line width=0.35pt] (4.317,-1.807) circle (0.20);
\fill[atomA, draw=atomLine, line width=0.35pt] (4.439,-0.807) circle (0.20);
\fill[atomA, draw=atomLine, line width=0.35pt] (4.561,0.193) circle (0.20);
\fill[atomA, draw=atomLine, line width=0.35pt] (4.683,1.193) circle (0.20);
\fill[atomA, draw=atomLine, line width=0.35pt] (4.803,2.192) circle (0.20);
\fill[atomA, draw=atomLine, line width=0.35pt] (4.925,3.189) circle (0.20);
\fill[atomA, draw=atomLine, line width=0.35pt] (5.047,4.186) circle (0.20);
\fill[atomA, draw=atomLine, line width=0.35pt] (5.169,5.184) circle (0.20);
\fill[atomA, draw=atomLine, line width=0.35pt] (5.289,6.180) circle (0.20);
\fill[atomA, draw=atomLine, line width=0.35pt] (5.410,7.174) circle (0.20);
\fill[atomA, draw=atomLine, line width=0.35pt] (4.716,-6.922) circle (0.20);
\fill[atomA, draw=atomLine, line width=0.35pt] (4.837,-5.919) circle (0.20);
\fill[atomA, draw=atomLine, line width=0.35pt] (4.957,-4.916) circle (0.20);
\fill[atomA, draw=atomLine, line width=0.35pt] (5.078,-3.913) circle (0.20);
\fill[atomA, draw=atomLine, line width=0.35pt] (5.199,-2.911) circle (0.20);
\fill[atomA, draw=atomLine, line width=0.35pt] (5.319,-1.910) circle (0.20);
\fill[atomA, draw=atomLine, line width=0.35pt] (5.440,-0.910) circle (0.20);
\fill[atomA, draw=atomLine, line width=0.35pt] (5.560,0.090) circle (0.20);
\fill[atomA, draw=atomLine, line width=0.35pt] (5.681,1.090) circle (0.20);
\fill[atomA, draw=atomLine, line width=0.35pt] (5.801,2.089) circle (0.20);
\fill[atomA, draw=atomLine, line width=0.35pt] (5.922,3.087) circle (0.20);
\fill[atomA, draw=atomLine, line width=0.35pt] (6.043,4.084) circle (0.20);
\fill[atomA, draw=atomLine, line width=0.35pt] (6.163,5.081) circle (0.20);
\fill[atomA, draw=atomLine, line width=0.35pt] (6.284,6.078) circle (0.20);
\fill[atomA, draw=atomLine, line width=0.35pt] (6.404,7.073) circle (0.20);
\fill[atomA, draw=atomLine, line width=0.35pt] (5.722,-7.020) circle (0.20);
\fill[atomA, draw=atomLine, line width=0.35pt] (5.842,-6.016) circle (0.20);
\fill[atomA, draw=atomLine, line width=0.35pt] (5.961,-5.013) circle (0.20);
\fill[atomA, draw=atomLine, line width=0.35pt] (6.081,-4.011) circle (0.20);
\fill[atomA, draw=atomLine, line width=0.35pt] (6.201,-3.010) circle (0.20);
\fill[atomA, draw=atomLine, line width=0.35pt] (6.320,-2.008) circle (0.20);
\fill[atomA, draw=atomLine, line width=0.35pt] (6.440,-1.008) circle (0.20);
\fill[atomA, draw=atomLine, line width=0.35pt] (6.560,-0.008) circle (0.20);
\fill[atomA, draw=atomLine, line width=0.35pt] (6.680,0.992) circle (0.20);
\fill[atomA, draw=atomLine, line width=0.35pt] (6.799,1.990) circle (0.20);
\fill[atomA, draw=atomLine, line width=0.35pt] (6.919,2.989) circle (0.20);
\fill[atomA, draw=atomLine, line width=0.35pt] (7.039,3.987) circle (0.20);
\fill[atomA, draw=atomLine, line width=0.35pt] (7.158,4.984) circle (0.20);
\fill[atomA, draw=atomLine, line width=0.35pt] (7.278,5.980) circle (0.20);
\fill[atomA, draw=atomLine, line width=0.35pt] (6.728,-7.113) circle (0.20);
\fill[atomA, draw=atomLine, line width=0.35pt] (6.847,-6.110) circle (0.20);
\fill[atomA, draw=atomLine, line width=0.35pt] (6.965,-5.107) circle (0.20);
\fill[atomA, draw=atomLine, line width=0.35pt] (7.084,-4.105) circle (0.20);
\fill[atomA, draw=atomLine, line width=0.35pt] (7.203,-3.104) circle (0.20);
\end{scope}
\dislT[fmRed]{(0.000,-3.055)}{0}
\dislT[fmRed]{(0.000,0.945)}{0}
\dislT[fmRed]{(0.000,4.945)}{0}
\coordinate (TB0) at (0,0.995);
\coordinate (TB1) at (0,4.995);

  \draw[dimline, fmRed] ($(TB0)+(1.3,0)$) -- node[right, slbl, fmRed]{$D$} ($(TB1)+(1.3,0)$);
\end{tikzpicture}
\caption{Симметричная граница наклона в простой кубической решётке~—
стенка краевых дислокаций через $D=4b$; угол разориентации
$\theta=b/D\approx14^\circ$ (положения атомов рассчитаны наложением упругих
полей дислокаций)}
\label{fig:lab-wall}
\end{figure}

\begin{primer}[Расстояние между дислокациями в субгранице]
В алюминии ($a=0{,}405$~нм, $b=a/\sqrt2=0{,}286$~нм) субграница наклона
с $\theta=2^\circ=0{,}035$~рад состоит из дислокаций, отстоящих на
$D=b/\theta=0{,}286/0{,}035\approx8{,}2$~нм, т.\,е. примерно через $29b$.
\end{primer}

В \emph{несимметричной} границе наклона (плоскость границы не делит угол
между решётками пополам) присутствуют дислокации с двумя разными
векторами Бюргерса. Малоугловая \term{граница кручения}~— сетка из двух
семейств параллельных \emph{винтовых} дислокаций; формула~\eqref{eq:wall-spacing}
справедлива и для неё, если $D$~— расстояние между дислокациями одного
семейства.

В общем случае дислокационное строение малоугловой границы задаёт
\term{формула Франка}:
\begin{equation}\label{eq:frank}
\vb d=\sum_i\vb b_i=2\sin\frac\theta2\,(\vb r\times\vb u),
\end{equation}
где $\vb r$~— произвольный вектор в плоскости границы, $\vb u$~— единичный
вектор оси разворота, $\vb d$~— суммарный вектор Бюргерса дислокаций,
которые пересекает вектор $\vb r$. Для симметричной границы наклона
($\vb r\perp\vb u$, оба в плоскости границы) $|\vb d|=2|\vb r|\sin(\theta/2)\approx|\vb r|\theta$:
вектор длины $|\vb r|$ пересекает $|\vb r|/D$ дислокаций, откуда снова
$D=b/\theta$.

\paragraph{Энергия малоугловой границы}
Энергия единицы площади границы растёт с углом разориентации
(\term{формула Рида--Шокли}):
\begin{keyeq}\label{eq:read-shockley}
E_\text{гр}=E_0\,\theta\,(A-\ln\theta),
\end{keyeq}
где $E_0\propto Gb$ ($G$~— модуль сдвига) и $A$~— константа, зависящая от
энергии ядра дислокации. Множитель $\theta$~— это плотность дислокаций в
стенке ($1/D=\theta/b$), а $-\ln\theta$ отражает взаимное экранирование
полей соседних дислокаций: чем они ближе, тем меньше энергия каждой.
Формула~\eqref{eq:read-shockley} имеет максимум при $\theta_m=e^{A-1}$, и её
удобно записать в приведённом виде
$E/E_m=(\theta/\theta_m)\,[1-\ln(\theta/\theta_m)]$ (рис.~\ref{fig:gb-energy},\,\textit{а}).
Практически она применима до $\theta\approx10$--$15^\circ$; выше энергия
границы почти не зависит от угла~— кроме особых разориентаций.

\subsection{Высокоугловые границы}\label{ssec:high-angle}

\paragraph{От аморфной прослойки к решёткам совпадения}
Первоначально границу зерна представляли \emph{аморфной прослойкой},
обеспечивающей сцепление зёрен. Эта модель противоречит опыту:
\begin{enumerate}
\item свойства аморфной границы не зависели бы от взаимной ориентации
зёрен, а на самом деле энергия границы, скорость диффузии по ней,
скорость миграции и зернограничное проскальзывание от разориентации
зависят; значит, граница имеет определённое строение;
\item чтобы зёрна неправильной формы могли проскальзывать, прослойка
должна быть толщиной порядка сотни атомных диаметров, а прямые
наблюдения (ионный проектор, электронная микроскопия высокого разрешения)
дают толщину всего в несколько межатомных расстояний.
\end{enumerate}
Затем предлагались модель \emph{переходной решётки} (атомы занимают
промежуточные положения между узлами решёток соседних зёрен) и
\emph{островковая модель Мотта}: граница состоит из участков «хорошего» и
«плохого» сопряжения решёток. В исходном виде модель Мотта не используют,
но сама идея чередования участков хорошего и плохого сопряжения лежит в
основе большинства современных моделей.

\paragraph{Решётка совпадающих узлов}
Мысленно продолжим решётки обоих зёрен на всё пространство. При
определённых, строго фиксированных осях и углах разворота часть узлов
одной решётки совпадает с узлами другой, и совпадающие узлы сами образуют
решётку (рис.~\ref{fig:csl5}).

\begin{opredelenie}
\term{Решётка совпадающих узлов} (РСУ, англ. CSL~— coincidence site lattice)~—
общая часть (пересечение) решёток двух зёрен, $L_\text{РСУ}=L_1\cap L_2$.
Величина $\Sigma$~— \term{обратная плотность совпадающих узлов}: число узлов
решётки, приходящееся на один совпадающий узел (отношение объёмов
элементарных ячеек РСУ и решётки).
\end{opredelenie}

Так, при повороте вокруг $[001]$ на угол
$\theta=2\operatorname{arctg}(1/3)=36{,}87^\circ$\ispr{В оригинале
$\theta=2\operatorname{arctg}(1/2)=36{,}9^\circ$. Но $2\operatorname{arctg}(1/2)=53{,}13^\circ$;
этот угол даёт ту же разориентацию $\Sigma5$ (так как $53{,}13^\circ=90^\circ-36{,}87^\circ$,
а поворот на $90^\circ$ вокруг $[001]$~— операция симметрии куба), но
значению $36{,}9^\circ$ соответствует $2\operatorname{arctg}(1/3)$.}
совпадает каждый пятый узел: $\Sigma=5$, а ячейка РСУ~— квадрат со
стороной $a\sqrt5$ (в трёх измерениях~— тетрагональная ячейка).
$\Sigma=1$ означает полное совпадение решёток (границы нет); $\Sigma=3$~—
двойник в ГЦК-решётке (поворот на $60^\circ$ вокруг $\nfam{111}$). Для
кубических решёток $\Sigma$ всегда \emph{нечётно}\ispr{В оригинале:
«число $\Sigma$ всегда простое». Это неверно: бывают $\Sigma=9$, 15, 21,
25, 27\ldots\ (табл.~\ref{tab:csl}). Верно другое: в кубических решётках
$\Sigma$ всегда нечётно.}.

\begin{figure}[ht]
\centering
\begin{tikzpicture}[scale=0.85]
  % создано tools/gen_boundaries.py: решётка совпадающих узлов Σ5 [001]
\begin{scope}
\clip (-3.25,-3.25) rectangle (3.25,3.25);
\draw[fmGreen!70!black, line width=0.5pt] (3,-9) -- ++(1,2) -- ++(-2,1);
\draw[fmGreen!70!black, line width=0.5pt] (1,-8) -- ++(1,2) -- ++(-2,1);
\draw[fmGreen!70!black, line width=0.5pt] (-1,-7) -- ++(1,2) -- ++(-2,1);
\draw[fmGreen!70!black, line width=0.5pt] (-3,-6) -- ++(1,2) -- ++(-2,1);
\draw[fmGreen!70!black, line width=0.5pt] (-5,-5) -- ++(1,2) -- ++(-2,1);
\draw[fmGreen!70!black, line width=0.5pt] (-7,-4) -- ++(1,2) -- ++(-2,1);
\draw[fmGreen!70!black, line width=0.5pt] (-9,-3) -- ++(1,2) -- ++(-2,1);
\draw[fmGreen!70!black, line width=0.5pt] (4,-7) -- ++(1,2) -- ++(-2,1);
\draw[fmGreen!70!black, line width=0.5pt] (2,-6) -- ++(1,2) -- ++(-2,1);
\draw[fmGreen!70!black, line width=0.5pt] (0,-5) -- ++(1,2) -- ++(-2,1);
\draw[fmGreen!70!black, line width=0.5pt] (-2,-4) -- ++(1,2) -- ++(-2,1);
\draw[fmGreen!70!black, line width=0.5pt] (-4,-3) -- ++(1,2) -- ++(-2,1);
\draw[fmGreen!70!black, line width=0.5pt] (-6,-2) -- ++(1,2) -- ++(-2,1);
\draw[fmGreen!70!black, line width=0.5pt] (-8,-1) -- ++(1,2) -- ++(-2,1);
\draw[fmGreen!70!black, line width=0.5pt] (5,-5) -- ++(1,2) -- ++(-2,1);
\draw[fmGreen!70!black, line width=0.5pt] (3,-4) -- ++(1,2) -- ++(-2,1);
\draw[fmGreen!70!black, line width=0.5pt] (1,-3) -- ++(1,2) -- ++(-2,1);
\draw[fmGreen!70!black, line width=0.5pt] (-1,-2) -- ++(1,2) -- ++(-2,1);
\draw[fmGreen!70!black, line width=0.5pt] (-3,-1) -- ++(1,2) -- ++(-2,1);
\draw[fmGreen!70!black, line width=0.5pt] (-5,0) -- ++(1,2) -- ++(-2,1);
\draw[fmGreen!70!black, line width=0.5pt] (-7,1) -- ++(1,2) -- ++(-2,1);
\draw[fmGreen!70!black, line width=0.5pt] (6,-3) -- ++(1,2) -- ++(-2,1);
\draw[fmGreen!70!black, line width=0.5pt] (4,-2) -- ++(1,2) -- ++(-2,1);
\draw[fmGreen!70!black, line width=0.5pt] (2,-1) -- ++(1,2) -- ++(-2,1);
\draw[fmGreen!70!black, line width=0.5pt] (0,0) -- ++(1,2) -- ++(-2,1);
\draw[fmGreen!70!black, line width=0.5pt] (-2,1) -- ++(1,2) -- ++(-2,1);
\draw[fmGreen!70!black, line width=0.5pt] (-4,2) -- ++(1,2) -- ++(-2,1);
\draw[fmGreen!70!black, line width=0.5pt] (-6,3) -- ++(1,2) -- ++(-2,1);
\draw[fmGreen!70!black, line width=0.5pt] (7,-1) -- ++(1,2) -- ++(-2,1);
\draw[fmGreen!70!black, line width=0.5pt] (5,0) -- ++(1,2) -- ++(-2,1);
\draw[fmGreen!70!black, line width=0.5pt] (3,1) -- ++(1,2) -- ++(-2,1);
\draw[fmGreen!70!black, line width=0.5pt] (1,2) -- ++(1,2) -- ++(-2,1);
\draw[fmGreen!70!black, line width=0.5pt] (-1,3) -- ++(1,2) -- ++(-2,1);
\draw[fmGreen!70!black, line width=0.5pt] (-3,4) -- ++(1,2) -- ++(-2,1);
\draw[fmGreen!70!black, line width=0.5pt] (-5,5) -- ++(1,2) -- ++(-2,1);
\draw[fmGreen!70!black, line width=0.5pt] (8,1) -- ++(1,2) -- ++(-2,1);
\draw[fmGreen!70!black, line width=0.5pt] (6,2) -- ++(1,2) -- ++(-2,1);
\draw[fmGreen!70!black, line width=0.5pt] (4,3) -- ++(1,2) -- ++(-2,1);
\draw[fmGreen!70!black, line width=0.5pt] (2,4) -- ++(1,2) -- ++(-2,1);
\draw[fmGreen!70!black, line width=0.5pt] (0,5) -- ++(1,2) -- ++(-2,1);
\draw[fmGreen!70!black, line width=0.5pt] (-2,6) -- ++(1,2) -- ++(-2,1);
\draw[fmGreen!70!black, line width=0.5pt] (-4,7) -- ++(1,2) -- ++(-2,1);
\draw[fmGreen!70!black, line width=0.5pt] (9,3) -- ++(1,2) -- ++(-2,1);
\draw[fmGreen!70!black, line width=0.5pt] (7,4) -- ++(1,2) -- ++(-2,1);
\draw[fmGreen!70!black, line width=0.5pt] (5,5) -- ++(1,2) -- ++(-2,1);
\draw[fmGreen!70!black, line width=0.5pt] (3,6) -- ++(1,2) -- ++(-2,1);
\draw[fmGreen!70!black, line width=0.5pt] (1,7) -- ++(1,2) -- ++(-2,1);
\draw[fmGreen!70!black, line width=0.5pt] (-1,8) -- ++(1,2) -- ++(-2,1);
\draw[fmGreen!70!black, line width=0.5pt] (-3,9) -- ++(1,2) -- ++(-2,1);
\fill[fmGreen!15] (0,0) -- (1,2) -- (-1,3) -- (-2,1) -- cycle;
\draw[fmGreen!70!black, line width=1.1pt] (0,0) -- (1,2) -- (-1,3) -- (-2,1) -- cycle;
\fill[fmGreen!60!black] (0.000,0.000) circle (0.045);
\fill[fmGreen!60!black] (-0.400,0.200) circle (0.045);
\fill[fmGreen!60!black] (-0.800,0.400) circle (0.045);
\fill[fmGreen!60!black] (-1.200,0.600) circle (0.045);
\fill[fmGreen!60!black] (-1.600,0.800) circle (0.045);
\fill[fmGreen!60!black] (-2.000,1.000) circle (0.045);
\fill[fmGreen!60!black] (0.200,0.400) circle (0.045);
\fill[fmGreen!60!black] (-0.200,0.600) circle (0.045);
\fill[fmGreen!60!black] (-0.600,0.800) circle (0.045);
\fill[fmGreen!60!black] (-1.000,1.000) circle (0.045);
\fill[fmGreen!60!black] (-1.400,1.200) circle (0.045);
\fill[fmGreen!60!black] (-1.800,1.400) circle (0.045);
\fill[fmGreen!60!black] (0.400,0.800) circle (0.045);
\fill[fmGreen!60!black] (0.000,1.000) circle (0.045);
\fill[fmGreen!60!black] (-0.400,1.200) circle (0.045);
\fill[fmGreen!60!black] (-0.800,1.400) circle (0.045);
\fill[fmGreen!60!black] (-1.200,1.600) circle (0.045);
\fill[fmGreen!60!black] (-1.600,1.800) circle (0.045);
\fill[fmGreen!60!black] (0.600,1.200) circle (0.045);
\fill[fmGreen!60!black] (0.200,1.400) circle (0.045);
\fill[fmGreen!60!black] (-0.200,1.600) circle (0.045);
\fill[fmGreen!60!black] (-0.600,1.800) circle (0.045);
\fill[fmGreen!60!black] (-1.000,2.000) circle (0.045);
\fill[fmGreen!60!black] (-1.400,2.200) circle (0.045);
\fill[fmGreen!60!black] (0.800,1.600) circle (0.045);
\fill[fmGreen!60!black] (0.400,1.800) circle (0.045);
\fill[fmGreen!60!black] (0.000,2.000) circle (0.045);
\fill[fmGreen!60!black] (-0.400,2.200) circle (0.045);
\fill[fmGreen!60!black] (-0.800,2.400) circle (0.045);
\fill[fmGreen!60!black] (-1.200,2.600) circle (0.045);
\fill[fmGreen!60!black] (1.000,2.000) circle (0.045);
\fill[fmGreen!60!black] (0.600,2.200) circle (0.045);
\fill[fmGreen!60!black] (0.200,2.400) circle (0.045);
\fill[fmGreen!60!black] (-0.200,2.600) circle (0.045);
\fill[fmGreen!60!black] (-0.600,2.800) circle (0.045);
\fill[fmGreen!60!black] (-1.000,3.000) circle (0.045);
\draw[fmRed, line width=0.7pt] (-2.6000,-3.2000) circle (0.15);
\draw[fmRed, line width=0.7pt] (-3.2000,-2.4000) circle (0.15);
\fill[fmBlue] (-3,-3) circle (0.09);
\fill[fmBlue] (-3,-2) circle (0.09);
\draw[fmRed, line width=0.7pt] (-1.2000,-3.4000) circle (0.15);
\fill[fmBlue] (-3,-1) circle (0.09);
\draw[fmRed, line width=0.7pt] (-1.8000,-2.6000) circle (0.15);
\fill[fmBlue] (-3,0) circle (0.09);
\draw[fmRed, line width=0.7pt] (-2.4000,-1.8000) circle (0.15);
\fill[fmBlue] (-3,1) circle (0.09);
\draw[fmRed, line width=0.7pt] (-3.0000,-1.0000) circle (0.15);
\draw[black, line width=1.1pt] (-3.0000,-1.0000) circle (0.25);
\fill[fmBlue] (-3,2) circle (0.09);
\fill[fmBlue] (-3,3) circle (0.09);
\fill[fmBlue] (-2,-3) circle (0.09);
\fill[fmBlue] (-2,-2) circle (0.09);
\draw[fmRed, line width=0.7pt] (-0.4000,-2.8000) circle (0.15);
\fill[fmBlue] (-2,-1) circle (0.09);
\draw[fmRed, line width=0.7pt] (-1.0000,-2.0000) circle (0.15);
\draw[black, line width=1.1pt] (-1.0000,-2.0000) circle (0.25);
\fill[fmBlue] (-2,0) circle (0.09);
\draw[fmRed, line width=0.7pt] (-1.6000,-1.2000) circle (0.15);
\fill[fmBlue] (-2,1) circle (0.09);
\draw[fmRed, line width=0.7pt] (-2.2000,-0.4000) circle (0.15);
\fill[fmBlue] (-2,2) circle (0.09);
\draw[fmRed, line width=0.7pt] (-2.8000,0.4000) circle (0.15);
\fill[fmBlue] (-2,3) circle (0.09);
\draw[fmRed, line width=0.7pt] (-3.4000,1.2000) circle (0.15);
\fill[fmBlue] (-1,-3) circle (0.09);
\draw[fmRed, line width=0.7pt] (1.0000,-3.0000) circle (0.15);
\draw[black, line width=1.1pt] (1.0000,-3.0000) circle (0.25);
\fill[fmBlue] (-1,-2) circle (0.09);
\draw[fmRed, line width=0.7pt] (0.4000,-2.2000) circle (0.15);
\fill[fmBlue] (-1,-1) circle (0.09);
\draw[fmRed, line width=0.7pt] (-0.2000,-1.4000) circle (0.15);
\fill[fmBlue] (-1,0) circle (0.09);
\draw[fmRed, line width=0.7pt] (-0.8000,-0.6000) circle (0.15);
\fill[fmBlue] (-1,1) circle (0.09);
\draw[fmRed, line width=0.7pt] (-1.4000,0.2000) circle (0.15);
\fill[fmBlue] (-1,2) circle (0.09);
\draw[fmRed, line width=0.7pt] (-2.0000,1.0000) circle (0.15);
\draw[black, line width=1.1pt] (-2.0000,1.0000) circle (0.25);
\fill[fmBlue] (-1,3) circle (0.09);
\draw[fmRed, line width=0.7pt] (-2.6000,1.8000) circle (0.15);
\draw[fmRed, line width=0.7pt] (-3.2000,2.6000) circle (0.15);
\draw[fmRed, line width=0.7pt] (2.4000,-3.2000) circle (0.15);
\fill[fmBlue] (0,-3) circle (0.09);
\draw[fmRed, line width=0.7pt] (1.8000,-2.4000) circle (0.15);
\fill[fmBlue] (0,-2) circle (0.09);
\draw[fmRed, line width=0.7pt] (1.2000,-1.6000) circle (0.15);
\fill[fmBlue] (0,-1) circle (0.09);
\draw[fmRed, line width=0.7pt] (0.6000,-0.8000) circle (0.15);
\fill[fmBlue] (0,0) circle (0.09);
\draw[fmRed, line width=0.7pt] (0.0000,0.0000) circle (0.15);
\draw[black, line width=1.1pt] (0.0000,0.0000) circle (0.25);
\fill[fmBlue] (0,1) circle (0.09);
\draw[fmRed, line width=0.7pt] (-0.6000,0.8000) circle (0.15);
\fill[fmBlue] (0,2) circle (0.09);
\draw[fmRed, line width=0.7pt] (-1.2000,1.6000) circle (0.15);
\fill[fmBlue] (0,3) circle (0.09);
\draw[fmRed, line width=0.7pt] (-1.8000,2.4000) circle (0.15);
\draw[fmRed, line width=0.7pt] (-2.4000,3.2000) circle (0.15);
\draw[fmRed, line width=0.7pt] (3.2000,-2.6000) circle (0.15);
\fill[fmBlue] (1,-3) circle (0.09);
\draw[fmRed, line width=0.7pt] (2.6000,-1.8000) circle (0.15);
\fill[fmBlue] (1,-2) circle (0.09);
\draw[fmRed, line width=0.7pt] (2.0000,-1.0000) circle (0.15);
\draw[black, line width=1.1pt] (2.0000,-1.0000) circle (0.25);
\fill[fmBlue] (1,-1) circle (0.09);
\draw[fmRed, line width=0.7pt] (1.4000,-0.2000) circle (0.15);
\fill[fmBlue] (1,0) circle (0.09);
\draw[fmRed, line width=0.7pt] (0.8000,0.6000) circle (0.15);
\fill[fmBlue] (1,1) circle (0.09);
\draw[fmRed, line width=0.7pt] (0.2000,1.4000) circle (0.15);
\fill[fmBlue] (1,2) circle (0.09);
\draw[fmRed, line width=0.7pt] (-0.4000,2.2000) circle (0.15);
\fill[fmBlue] (1,3) circle (0.09);
\draw[fmRed, line width=0.7pt] (-1.0000,3.0000) circle (0.15);
\draw[black, line width=1.1pt] (-1.0000,3.0000) circle (0.25);
\fill[fmBlue] (2,-3) circle (0.09);
\draw[fmRed, line width=0.7pt] (3.4000,-1.2000) circle (0.15);
\fill[fmBlue] (2,-2) circle (0.09);
\draw[fmRed, line width=0.7pt] (2.8000,-0.4000) circle (0.15);
\fill[fmBlue] (2,-1) circle (0.09);
\draw[fmRed, line width=0.7pt] (2.2000,0.4000) circle (0.15);
\fill[fmBlue] (2,0) circle (0.09);
\draw[fmRed, line width=0.7pt] (1.6000,1.2000) circle (0.15);
\fill[fmBlue] (2,1) circle (0.09);
\draw[fmRed, line width=0.7pt] (1.0000,2.0000) circle (0.15);
\draw[black, line width=1.1pt] (1.0000,2.0000) circle (0.25);
\fill[fmBlue] (2,2) circle (0.09);
\draw[fmRed, line width=0.7pt] (0.4000,2.8000) circle (0.15);
\fill[fmBlue] (2,3) circle (0.09);
\fill[fmBlue] (3,-3) circle (0.09);
\fill[fmBlue] (3,-2) circle (0.09);
\fill[fmBlue] (3,-1) circle (0.09);
\draw[fmRed, line width=0.7pt] (3.0000,1.0000) circle (0.15);
\draw[black, line width=1.1pt] (3.0000,1.0000) circle (0.25);
\fill[fmBlue] (3,0) circle (0.09);
\draw[fmRed, line width=0.7pt] (2.4000,1.8000) circle (0.15);
\fill[fmBlue] (3,1) circle (0.09);
\draw[fmRed, line width=0.7pt] (1.8000,2.6000) circle (0.15);
\fill[fmBlue] (3,2) circle (0.09);
\draw[fmRed, line width=0.7pt] (1.2000,3.4000) circle (0.15);
\fill[fmBlue] (3,3) circle (0.09);
\draw[fmRed, line width=0.7pt] (3.2000,2.4000) circle (0.15);
\draw[fmRed, line width=0.7pt] (2.6000,3.2000) circle (0.15);
\end{scope}
\draw[atomLine] (-3.25,-3.25) rectangle (3.25,3.25);
\begin{scope}[shift={(7.6,0)}]
\clip (-3.25,-3.25) rectangle (3.25,3.25);
\fill[fmBlue!7] (-3.25,-3.25) rectangle (0,3.25);
\fill[fmRed!7] (0,-3.25) rectangle (3.25,3.25);
\node[atom=2.3mm] at (-3.4785,2.2136) {};
\node[atom=2.3mm] at (-3.1623,3.1623) {};
\node[atom=2.3mm] at (-3.4785,-0.9487) {};
\node[atom=2.3mm] at (-3.1623,0.0000) {};
\node[atom=2.3mm] at (-2.8460,0.9487) {};
\node[atom=2.3mm] at (-2.5298,1.8974) {};
\node[atom=2.3mm] at (-2.2136,2.8460) {};
\node[atom=2.3mm] at (-3.1623,-3.1623) {};
\node[atom=2.3mm] at (-2.8460,-2.2136) {};
\node[atom=2.3mm] at (-2.5298,-1.2649) {};
\node[atom=2.3mm] at (-2.2136,-0.3162) {};
\node[atom=2.3mm] at (-1.8974,0.6325) {};
\node[atom=2.3mm] at (-1.5811,1.5811) {};
\node[atom=2.3mm] at (-1.2649,2.5298) {};
\node[atom=2.3mm] at (-0.9487,3.4785) {};
\node[atom=2.3mm] at (-2.2136,-3.4785) {};
\node[atom=2.3mm] at (-1.8974,-2.5298) {};
\node[atom=2.3mm, fill=atomI] at (0.0000,-3.1623) {};
\draw[black, line width=1.1pt] (0.0000,-3.1623) circle (0.25);
\node[atom=2.3mm] at (-1.5811,-1.5811) {};
\node[atom=2.3mm] at (-1.2649,-0.6325) {};
\node[atom=2.3mm] at (-0.9487,0.3162) {};
\node[atom=2.3mm] at (-0.6325,1.2649) {};
\node[atom=2.3mm] at (-0.3162,2.2136) {};
\node[atom=2.3mm, fill=atomI] at (0.0000,3.1623) {};
\draw[black, line width=1.1pt] (0.0000,3.1623) circle (0.25);
\node[atom=2.3mm] at (-0.9487,-2.8460) {};
\node[atomB=2.3mm] at (0.9487,-2.8460) {};
\node[atom=2.3mm] at (-0.6325,-1.8974) {};
\node[atomB=2.3mm] at (0.6325,-1.8974) {};
\node[atom=2.3mm] at (-0.3162,-0.9487) {};
\node[atomB=2.3mm] at (0.3162,-0.9487) {};
\node[atom=2.3mm, fill=atomI] at (0.0000,0.0000) {};
\draw[black, line width=1.1pt] (0.0000,0.0000) circle (0.25);
\node[atomB=2.3mm] at (2.2136,-3.4785) {};
\node[atomB=2.3mm] at (1.8974,-2.5298) {};
\node[atomB=2.3mm] at (1.5811,-1.5811) {};
\node[atomB=2.3mm] at (1.2649,-0.6325) {};
\node[atomB=2.3mm] at (0.9487,0.3162) {};
\node[atomB=2.3mm] at (0.6325,1.2649) {};
\node[atomB=2.3mm] at (0.3162,2.2136) {};
\node[atomB=2.3mm] at (3.1623,-3.1623) {};
\node[atomB=2.3mm] at (2.8460,-2.2136) {};
\node[atomB=2.3mm] at (2.5298,-1.2649) {};
\node[atomB=2.3mm] at (2.2136,-0.3162) {};
\node[atomB=2.3mm] at (1.8974,0.6325) {};
\node[atomB=2.3mm] at (1.5811,1.5811) {};
\node[atomB=2.3mm] at (1.2649,2.5298) {};
\node[atomB=2.3mm] at (0.9487,3.4785) {};
\node[atomB=2.3mm] at (3.4785,-0.9487) {};
\node[atomB=2.3mm] at (3.1623,0.0000) {};
\node[atomB=2.3mm] at (2.8460,0.9487) {};
\node[atomB=2.3mm] at (2.5298,1.8974) {};
\node[atomB=2.3mm] at (2.2136,2.8460) {};
\node[atomB=2.3mm] at (3.4785,2.2136) {};
\node[atomB=2.3mm] at (3.1623,3.1623) {};
\end{scope}
\draw[atomLine] (7.6-3.25,-3.25) rectangle (7.6+3.25,3.25);
\draw[fmGray, dashed, line width=0.6pt] (7.6,-3.25) -- (7.6,3.25);

  \node[font=\small] at (0,-3.85) {\textit{а}};
  \node[font=\small] at (7.6,-3.85) {\textit{б}};
\end{tikzpicture}
\caption{Решётка совпадающих узлов $\Sigma5$ (поворот на $36{,}87^\circ$
вокруг $[001]$, показана плоскость $(001)$). \textit{а}~— наложение двух
решёток: синие точки и красные кружки; совпадающие узлы обведены, ячейка
РСУ выделена зелёным; зелёные точки внутри неё~— узлы полной решётки
наложения (ПРН), шаг которой в $\sqrt5$ раз меньше периода. \textit{б}~—
симметричная граница наклона $\Sigma5\,(310)$: слева решётка первого
зерна, справа второго; на границе (штриховая линия) каждый узел~—
совпадающий (красные атомы)}
\label{fig:csl5}
\end{figure}

\begin{table}[ht]
\centering
\caption{Разориентации кубических решёток с малыми $\Sigma$ (ось и
наименьший угол поворота) и допустимые отклонения по критерию
Брэндона~\eqref{eq:brandon}}
\label{tab:csl}
\small
\begin{tabular}{@{}rclc@{\hspace{10mm}}rclc@{}}
\toprule
$\Sigma$ & Ось & Угол & $\Delta\theta_\text{max}$ & $\Sigma$ & Ось & Угол & $\Delta\theta_\text{max}$\\
\midrule
3   & $\nfam{111}$ & $60{,}00^\circ$ & $8{,}7^\circ$ & 15  & $\nfam{210}$ & $48{,}19^\circ$ & $3{,}9^\circ$\\
5   & $\nfam{100}$ & $36{,}87^\circ$ & $6{,}7^\circ$ & 17a & $\nfam{100}$ & $28{,}07^\circ$ & $3{,}6^\circ$\\
7   & $\nfam{111}$ & $38{,}21^\circ$ & $5{,}7^\circ$ & 19a & $\nfam{110}$ & $26{,}53^\circ$ & $3{,}4^\circ$\\
9   & $\nfam{110}$ & $38{,}94^\circ$ & $5{,}0^\circ$ & 21a & $\nfam{111}$ & $21{,}79^\circ$ & $3{,}3^\circ$\\
11  & $\nfam{110}$ & $50{,}48^\circ$ & $4{,}5^\circ$ & 25a & $\nfam{100}$ & $16{,}26^\circ$ & $3{,}0^\circ$\\
13a & $\nfam{100}$ & $22{,}62^\circ$ & $4{,}2^\circ$ & 27a & $\nfam{110}$ & $31{,}59^\circ$ & $2{,}9^\circ$\\
13b & $\nfam{111}$ & $27{,}80^\circ$ & $4{,}2^\circ$ & 29a & $\nfam{100}$ & $43{,}60^\circ$ & $2{,}8^\circ$\\
\bottomrule
\end{tabular}
\end{table}

Для поворота вокруг оси $[hkl]$ угол, дающий РСУ, удовлетворяет условию
$\operatorname{tg}(\theta/2)=(n/m)\sqrt{N}$, где $N=h^2+k^2+l^2$, а $m$ и
$n$~— взаимно простые целые; тогда $\Sigma=m^2+Nn^2$, делённое на 2 столько
раз, сколько возможно (соотношения Кронберга--Вильсона в форме Ранганатана).
Например, для $\Sigma5$: ось $[100]$ ($N=1$), $m=3$, $n=1$:
$\Sigma=(9+1)/2=5$; для $\Sigma3$: ось $[111]$, $m=3$, $n=1$:
$(9+3)/4=3$ (табл.~\ref{tab:csl}).

\paragraph{Специальные границы}
Если граница совпадает с плотноупакованной плоскостью РСУ, в ней много
совпадающих узлов, решётки сопрягаются хорошо, и энергия границы мала. Такие
границы называют \term{специальными} (особыми); на зависимости энергии от
угла разворота им соответствуют резкие провалы (рис.~\ref{fig:gb-energy},\,\textit{б}).
Практически специальными считают границы с $1\le\Sigma\lesssim25$--$29$.
Самая «специальная»~— когерентная двойниковая граница $\Sigma3$ по $\plfam{111}$
в ГЦК-металлах: её энергия в десятки раз меньше, чем у обычной
высокоугловой границы (у меди $\approx0{,}02$ и $\approx0{,}6$~Дж/м$^2$).

\begin{figure}[ht]
\centering
\begin{tikzpicture}
\begin{axis}[name=rs, width=0.44\textwidth, height=0.36\textwidth,
  xlabel={$\theta/\theta_m$}, ylabel={$E/E_m$}, xmin=0, xmax=2.2, ymin=0, ymax=1.1,
  tick label style={font=\footnotesize}, label style={font=\small}]
  \addplot[fmBlue, line width=1.1pt, domain=0.001:1, samples=120] {x*(1-ln(x))};
  \addplot[fmBlue, line width=1.1pt, dashed, domain=1:2.2, samples=2] {1};
  \node[font=\footnotesize, anchor=north west, align=left] at (axis cs:0.35,0.55) {Рид--Шокли};
  \node[font=\footnotesize, anchor=north] at (axis cs:1.6,0.95) {высокоугловые};
\end{axis}
\begin{axis}[at={(rs.east)}, anchor=west, xshift=12mm, width=0.5\textwidth, height=0.36\textwidth,
  xlabel={угол разворота вокруг $\nfam{110}$, град}, ylabel={$E/E_\text{пл}$},
  xmin=0, xmax=180, ymin=0, ymax=1.15, xtick={0,30,...,180},
  tick label style={font=\footnotesize}, label style={font=\small}]
  \addplot[fmRed, line width=1.1pt] table[x=theta, y=E] {figures/gb-energy.dat};
  \foreach \t/\l/\y in {38.94/9/0.78, 50.48/11/0.53, 109.47/3/0.53, 129.52/11/0.73, 141.06/9/0.78}{
    \edef\tmp{\noexpand\node[font=\noexpand\scriptsize, anchor=north] at (axis cs:\t,\y) {$\noexpand\Sigma\l$};}\tmp}
  \node[font=\scriptsize, anchor=west] at (axis cs:73,0.07) {$\Sigma3$};
\end{axis}
\end{tikzpicture}
\caption{Энергия границы зерна: \textit{а}~— малоугловые границы,
формула Рида--Шокли~\eqref{eq:read-shockley} в приведённых координатах;
\textit{б}~— схема зависимости энергии симметричных границ наклона вокруг
$\nfam{110}$ в ГЦК-металле от угла разворота: провалы при специальных
разориентациях; глубже всех~— когерентный двойник $\Sigma3$ $(111)$ при
$70{,}5^\circ$}
\label{fig:gb-energy}
\end{figure}

Если граница идёт под небольшим углом к плотноупакованной плоскости РСУ,
она становится \emph{ступенчатой} (фасетированной): большей частью своей
площади она располагается по плоскостям с максимальной плотностью
совпадающих узлов.

Специальные свойства сохраняются и при небольших отклонениях
разориентации от специальной. Отклонение «принимают на себя»
\term{зернограничные дислокации}: они приспосабливают (аккомодируют)
структуру границы к фактической разориентации~— так же, как решёточные
дислокации образуют малоугловую границу в совершенном кристалле.
Предельное отклонение, при котором это ещё возможно, оценивают по
\term{критерию Брэндона}:
\begin{keyeq}\label{eq:brandon}
\Delta\theta_\text{max}=\frac{\theta_0}{\sqrt\Sigma},\qquad\theta_0\approx15^\circ
\end{keyeq}
(табл.~\ref{tab:csl})\ispr{В оригинале (формула (1.152)) критерий
записан как $\Delta\theta=(10\div15)/(\pi\sqrt\Sigma)$~рад, что даёт для
$\Sigma3$ около $160^\circ$. Критерий Брэндона: $\Delta\theta=\theta_0/\sqrt\Sigma$,
где $\theta_0=10$--$15^\circ$ (граница малоугловых границ);
в радианах $\theta_0=(10\div15)\pi/180$.}. Чем больше $\Sigma$, тем реже
совпадающие узлы и тем меньше допустимое отклонение.

\paragraph{Пространство разориентаций}
Разворот двух зёрен задаётся осью $\vb l$ и углом $\theta$ (или матрицей
поворота $A_0(\vb l,\theta)$, формула~\eqref{eq:R-axis}). Все возможные
развороты можно изобразить точками шара радиуса $\pi$ (направление из центра~—
ось, расстояние~— угол). С учётом веса, как и для углов Эйлера, «объём»
этого пространства равен $8\pi^2$.

Но из-за симметрии кристаллов один и тот же разворот описывается
несколькими способами. Если $R_i$~— поворотные элементы симметрии
кристалла, то эквивалентны описания
\begin{equation}\label{eq:equiv-misor}
A_i=A_0R_i^{-1}
\end{equation}
(поворот кристалла как целого операцией симметрии ничего не меняет). Для
кубических кристаллов (24 поворотных элемента группы $432$) существует
\vazhno{24 эквивалентных описания} разориентации с разными осями и углами
$\theta_i$. Описание с наименьшим углом называют \term{разориентировкой}
(англ. disorientation); обычно именно его и приводят.

Поэтому угол разориентации не может быть сколь угодно большим.
\term{Предельный угол} вокруг оси $[mnp]$~— наибольший угол, при котором
описание ещё остаётся наименьшим; при $\theta>\theta_\text{пр}$ найдётся
эквивалентное описание с меньшим углом. Для кубических кристаллов
предельные углы равны $45^\circ$ вокруг $\nfam{100}$, $60{,}7^\circ$ вокруг
$\nfam{110}$ и $60^\circ$ вокруг $\nfam{111}$, а наибольший возможный угол
разориентации~— $62{,}8^\circ$ вокруг оси $\nfam{1\,1\,\sqrt2-1}$, близкой
к $\nfam{221}$. Если зёрна ориентированы случайно, углы разориентации
распределены по закону Маккензи (рис.~\ref{fig:mackenzie}): чаще всего
встречаются углы около $45^\circ$, средний угол $40{,}7^\circ$, а
малоугловых границ мало.

\begin{figure}[ht]
\centering
\begin{tikzpicture}
\begin{axis}[width=0.62\textwidth, height=0.34\textwidth,
  xlabel={угол разориентации $\theta$, град}, ylabel={доля, 1/град},
  xmin=0, xmax=65, ymin=0, ymax=0.048, xtick={0,10,...,60},
  yticklabel style={/pgf/number format/fixed, /pgf/number format/precision=2},
  scaled y ticks=false,
  tick label style={font=\footnotesize}, label style={font=\small}]
  \addplot[const plot mark mid, fill=fmBlue!25, draw=fmBlue!70] table[x=theta, y=p] {figures/gb-mackenzie.dat} \closedcycle;
  \draw[fmRed, dashed] (axis cs:62.8,0) -- (axis cs:62.8,0.044);
  \node[font=\scriptsize, fmRed, anchor=south east] at (axis cs:62.8,0.040) {$62{,}8^\circ$};
\end{axis}
\end{tikzpicture}
\caption{Распределение Маккензи: доля границ с данным углом разориентации
в поликристалле со случайно ориентированными кубическими зёрнами (расчёт
методом Монте-Карло; учтены 24 эквивалентных описания~\eqref{eq:equiv-misor})}
\label{fig:mackenzie}
\end{figure}

\paragraph{Полная решётка наложения и O-решётка}
Если РСУ~— пересечение решёток двух зёрен, то \term{полная решётка наложения}
(ПРН, англ. DSC~— displacement shift complete lattice)~— наименьшая
решётка, содержащая узлы \emph{обеих} решёток\ispr{В оригинале ПРН
определена как объединение множеств узлов $L_1\cup L_2$. Объединение двух
решёток, вообще говоря, не является решёткой; ПРН~— наименее плотная
решётка, содержащая это объединение (порождённая им).}. Её векторы~—
наименьшие смещения одной решётки относительно другой, сохраняющие картину
совпадений (рис.~\ref{fig:csl5},\,\textit{а}). Объём ячейки РСУ в $\Sigma$ раз
больше, а объём ячейки ПРН в $\Sigma$ раз меньше объёма ячейки решётки
$V$:
\[
V_\text{РСУ}=\Sigma V,\qquad V_\text{ПРН}=V/\Sigma,\qquad
V_\text{РСУ}V_\text{ПРН}=V^2 .
\]
Векторы ПРН~— это возможные векторы Бюргерса зернограничных дислокаций;
они \emph{короче} векторов решётки (для $\Sigma5$ в плоскости $(001)$~—
в $\sqrt5$ раз). Наиболее плотноупакованное направление ПРН задаёт вектор
Бюргерса этих дислокаций.

Для расчёта дислокационного строения \emph{любой} границы (а не только
специальной) Боллман ввёл \term{O-решётку}~— множество точек, в которых
совпадают не только узлы, но и любые физически эквивалентные точки
элементарных ячеек двух решёток. O-решётка существует при любом развороте;
РСУ~— её частный случай. Для малоугловых границ шаг O-решётки
$b/(2\sin\frac\theta2)$ совпадает с расстоянием между дислокациями~—
формула Франка получается как частный случай.

\begin{itogi}
\begin{itemize}
\item Граница задаётся 5 макроскопическими параметрами (3~— разворот, 2~—
нормаль) и 3 микроскопическими (сдвиг $\vb t$).
\item Граница наклона: ось разворота в плоскости границы; кручения:
ось перпендикулярна границе.
\item Малоугловые границы~— стенки и сетки дислокаций: $D=b/\theta$;
формула Франка $\vb d=2\sin\frac\theta2(\vb r\times\vb u)$; энергия
$E=E_0\theta(A-\ln\theta)$.
\item Высокоугловые границы описывают РСУ: $\Sigma$~— число узлов на один
совпадающий (в кубических решётках нечётно); специальные границы
($\Sigma\lesssim25$--$29$) имеют пониженную энергию; критерий Брэндона
$\Delta\theta=15^\circ/\sqrt\Sigma$.
\item У кубических кристаллов 24 эквивалентных описания разворота;
наибольший угол разориентации $62{,}8^\circ$.
\item ПРН (DSC)~— решётка векторов Бюргерса зернограничных дислокаций,
$V_\text{РСУ}V_\text{ПРН}=V^2$; O-решётка Боллмана обобщает РСУ на любые
развороты.
\end{itemize}
\end{itogi}

\enlargethispage{2\baselineskip}
\begin{proverka}
\item Сколько параметров задают плоскую границу зерна с макроскопической
и с атомной точки зрения?
\item Чем граница наклона отличается от границы кручения? Из каких
дислокаций они состоят, если малоугловые?
\item Субграница наклона в меди ($b=0{,}256$~нм) содержит дислокации через
50~нм. Найдите угол разориентации.
\item Почему поворот на $53{,}13^\circ$ вокруг $[001]$ даёт ту же РСУ $\Sigma5$,
что и поворот на $36{,}87^\circ$?
\item Почему двойниковая граница $\Sigma3$ в ГЦК-металлах имеет очень
малую энергию?
\item Может ли угол разориентации двух кубических зёрен равняться
$70^\circ$?
\end{proverka}
